% SPDX-License-Identifier: MIT
% Copyright (c) 2026 Aymix — LaTeX design system, reusable under the MIT Licence.
% =====================================================================
%  THE BACKEND ENGINEER'S LAB NOTEBOOK  —  shared design system
%  Companion workbook to the "Java & Spring Boot 14-Day Speedrun Roadmap"
%  Engine: pdflatex   |   No shell-escape required (listings, not minted)
% =====================================================================

% ----- core -----
\usepackage[T1]{fontenc}
\usepackage[utf8]{inputenc}
\usepackage{lmodern}
\usepackage{amssymb}
\usepackage{etoolbox}
\usepackage[table,dvipsnames]{xcolor}
\usepackage{graphicx}
\usepackage{array}
\usepackage{booktabs}
\usepackage{tabularx}
\usepackage{multirow}
\usepackage{multicol}
\usepackage{enumitem}
\usepackage{ragged2e}
\usepackage[a4paper,top=2.5cm,bottom=2.3cm,left=2.25cm,right=2.25cm,headheight=17pt,headsep=6mm,footskip=12mm]{geometry}

% ----- fonts: Fira Sans body + Source Code Pro mono -----
\usepackage[sfdefault,scaled=1.0]{FiraSans}
\usepackage[scaled=0.84]{sourcecodepro}
\usepackage[defaultsans]{lato} % fallback safety (not default)
\renewcommand{\familydefault}{\sfdefault}
\usepackage{microtype}

% ----- fontawesome icons -----
\usepackage{fontawesome5}

% =====================================================================
%  COLOUR SYSTEM
% =====================================================================
\definecolor{cInk}{HTML}{1C2230}        % near-black body text
\definecolor{cBrand}{HTML}{243B6B}      % deep indigo — primary
\definecolor{cBrandDeep}{HTML}{17264A}  % darker indigo (headers)
\definecolor{cAccent}{HTML}{E08A1E}     % warm amber — notebook accent
\definecolor{cRule}{HTML}{C9CFDA}       % light rules / borders
\definecolor{cFigBg}{HTML}{FBFCFE}      % figure background
\definecolor{cCodeBg}{HTML}{1E2536}     % code background (dark)
\definecolor{cCodeFrame}{HTML}{2C3550}
\definecolor{cCodeTitle}{HTML}{111725}

% semantic callout colours
\definecolor{cSenior}{HTML}{0F766E}     % teal-green — senior mindset
\definecolor{cProd}{HTML}{1D6FB8}       % blue — production tip
\definecolor{cPerf}{HTML}{B45309}       % amber-brown — performance
\definecolor{cSecurity}{HTML}{B3261E}   % red — security warning
\definecolor{cDb}{HTML}{6D28D9}         % purple — database insight
\definecolor{cInterview}{HTML}{3730A3}  % indigo — interview insight
\definecolor{cMiscon}{HTML}{C2410C}     % orange — common misconception
\definecolor{cBehind}{HTML}{4C1D95}     % violet — behind the scenes
\definecolor{cGood}{HTML}{1E7A46}       % green — best practices
\definecolor{cBad}{HTML}{B3261E}        % red — common mistakes
\definecolor{cThink}{HTML}{7A5A12}      % ochre — thinking / reflection
\definecolor{cLab}{HTML}{0E7490}        % cyan-teal — labs
\definecolor{cDebug}{HTML}{9A3412}      % burnt orange — debugging
\definecolor{cBuild}{HTML}{155E75}      % deep cyan — build-along

% =====================================================================
%  CODE (listings inside tcolorbox)
% =====================================================================
\usepackage{listings}
\usepackage[most]{tcolorbox}
\tcbuselibrary{listings,skins,breakable}

\definecolor{lstKw}{HTML}{7CC5FF}
\definecolor{lstKw2}{HTML}{C792EA}
\definecolor{lstStr}{HTML}{ECC48D}
\definecolor{lstCom}{HTML}{7E8AA3}
\definecolor{lstNum}{HTML}{F78C6C}
\definecolor{lstTxt}{HTML}{E6EAF2}
\definecolor{lstAnno}{HTML}{82E0C6}

\lstset{
  basicstyle=\ttfamily\footnotesize\color{lstTxt},
  keywordstyle=\color{lstKw}\bfseries,
  keywordstyle=[2]\color{lstKw2},
  commentstyle=\color{lstCom}\itshape,
  stringstyle=\color{lstStr},
  numberstyle=\ttfamily\scriptsize\color{lstCom},
  showstringspaces=false,
  keepspaces=true,
  upquote=true,
  columns=fullflexible,
  breaklines=true,
  breakatwhitespace=false,
  tabsize=2,
  extendedchars=true,
  literate={→}{{$\rightarrow$}}1 {←}{{$\leftarrow$}}1 {≤}{{$\leq$}}1 {≥}{{$\geq$}}1
}

% Java (built-in) + annotation highlighting
\lstdefinestyle{stJava}{
  language=Java,
  morekeywords={var,record,sealed,yield,permits},
  moredelim=[is][\color{lstAnno}]{@@}{@@},
}
\lstdefinestyle{stSql}{language=SQL, morekeywords={ANALYZE,EXPLAIN,SERIALIZABLE,UNCOMMITTED,COMMITTED,REPEATABLE,ILIKE,RETURNING}}

\lstdefinelanguage{yamlx}{
  keywords={true,false,null},
  keywordstyle=\color{lstKw2},
  comment=[l]{\#},
  morestring=[b]",
  moredelim=[l][\color{lstKw}]{},
}
\lstdefinestyle{stYaml}{language=yamlx, basicstyle=\ttfamily\footnotesize\color{lstTxt}}

\lstdefinelanguage{dockerx}{
  morekeywords={FROM,AS,WORKDIR,COPY,RUN,CMD,ENTRYPOINT,EXPOSE,ENV,USER,ARG,VOLUME,LABEL,ADD,HEALTHCHECK},
  keywordstyle=\color{lstKw}\bfseries,
  comment=[l]{\#},
  morestring=[b]",
}
\lstdefinestyle{stDocker}{language=dockerx}

\lstdefinelanguage{bashx}{
  morekeywords={curl,docker,mvn,java,git,echo,export,cd,psql,redis-cli,kubectl,grep,cat,sudo,run,build,push,compose,up,down},
  keywordstyle=\color{lstKw}\bfseries,
  comment=[l]{\#},
  morestring=[b]",morestring=[b]',
}
\lstdefinestyle{stBash}{language=bashx}

\lstdefinelanguage{httpx}{
  morekeywords={GET,POST,PUT,PATCH,DELETE,HTTP,Authorization,Bearer,Content-Type,Accept,Location},
  keywordstyle=\color{lstKw}\bfseries,
  morestring=[b]",
}
\lstdefinestyle{stHttp}{language=httpx}

\lstdefinestyle{stProps}{
  morecomment=[l]{\#},
  morekeywords={true,false},
  keywordstyle=\color{lstKw2},
}

% base code box
\tcbset{
  wbcode/.style={
    enhanced, breakable,
    colback=cCodeBg, colframe=cCodeFrame, boxrule=0.6pt,
    arc=1.2mm, outer arc=1.2mm,
    left=3mm, right=3mm, top=1.8mm, bottom=1.8mm, boxsep=0.5mm,
    coltitle=white, fonttitle=\ttfamily\bfseries\scriptsize,
    colbacktitle=cCodeTitle,
    attach boxed title to top left={xshift=2.2mm, yshift=-2.4mm},
    boxed title style={colback=cCodeTitle, colframe=cCodeTitle, arc=0.6mm, boxrule=0pt},
    listing only,
    top=3.2mm,
  }
}
\newtcblisting{javabox}[1][Java]{wbcode, listing style=stJava, title={\ #1\ }}
\newtcblisting{sqlbox}[1][SQL]{wbcode, listing style=stSql, title={\ #1\ }}
\newtcblisting{yamlbox}[1][application.yml]{wbcode, listing style=stYaml, title={\ #1\ }}
\newtcblisting{dockerbox}[1][Dockerfile]{wbcode, listing style=stDocker, title={\ #1\ }}
\newtcblisting{bashbox}[1][shell]{wbcode, listing style=stBash, title={\ #1\ }}
\newtcblisting{httpbox}[1][HTTP]{wbcode, listing style=stHttp, title={\ #1\ }}
\newtcblisting{propsbox}[1][properties]{wbcode, listing style=stProps, title={\ #1\ }}

% inline code — colored monospace, line-break friendly (never overflows the margin)
\newcommand{\code}[1]{{\ttfamily\small\color{cBrandDeep}#1}}
% boxed chip — ONLY for very short labels/keys (<= ~14 chars); do not use for long identifiers
\newcommand{\chip}[1]{{\ttfamily\footnotesize\colorbox{cRule!35}{\hspace{0.3mm}\textcolor{cBrandDeep}{#1}\hspace{0.3mm}}}}
\newcommand{\anno}[1]{{\ttfamily\small\textbf{\color{cDb}#1}}}

% =====================================================================
%  CALLOUT BOXES
% =====================================================================
\tcbset{
  callout/.style={
    enhanced, breakable,
    colframe=black, colback=white, coltitle=white,
    fonttitle=\sffamily\bfseries\small,
    fontupper=\small,
    arc=1.4mm, outer arc=1.4mm,
    boxrule=0.5pt, leftrule=3.4pt,
    left=3.2mm, right=3.2mm, top=2.4mm, bottom=2.4mm, boxsep=0.8mm,
    toptitle=0.8mm, bottomtitle=0.8mm, lefttitle=1mm,
  }
}
% factory: #1 env, #2 colour, #3 icon+title
\newcommand{\definecallout}[3]{%
  \newtcolorbox{#1}[1][]{callout, colframe=#2, colback=#2!6,
    title={#3\ifblank{##1}{}{\ \textmd{\textbullet}\ ##1}}}%
}
\definecallout{seniornote}{cSenior}{\faLightbulb\enspace How senior engineers think}
\definecallout{prodtip}{cProd}{\faRocket\enspace Production tip}
\definecallout{perfnote}{cPerf}{\faBolt\enspace Performance note}
\definecallout{secwarn}{cSecurity}{\faUserShield\enspace Security warning}
\definecallout{dbinsight}{cDb}{\faDatabase\enspace Database insight}
\definecallout{interview}{cInterview}{\faUserTie\enspace Interview insight}
\definecallout{misconception}{cMiscon}{\faExclamationTriangle\enspace Common misconception}
\definecallout{behind}{cBehind}{\faEye\enspace Behind the scenes}
\definecallout{whyexists}{cBrand}{\faQuestionCircle\enspace Why does this exist?}
\definecallout{keyidea}{cBrand}{\faKey\enspace Key idea}

% Best practices / Mistakes — two-column friendly, no title bar variant
\newtcolorbox{bestpractices}{callout, colframe=cGood, colback=cGood!6,
  title={\faCheckCircle\enspace Best practices — what experienced engineers do}}
\newtcolorbox{mistakes}{callout, colframe=cBad, colback=cBad!6,
  title={\faTimesCircle\enspace Common mistakes beginners make}}

% =====================================================================
%  STRUCTURAL SECTION HEADERS  (11 canonical parts per chapter)
% =====================================================================
\usepackage{tikz}
\usetikzlibrary{arrows.meta,positioning,shapes.geometric,shapes.misc,fit,backgrounds,calc,chains}

\newcommand{\wbpart}[3][cBrand]{%
  \par\addvspace{6.5mm}%
  \noindent\begin{tikzpicture}[baseline=(t.base)]
    \node[fill=#1, text=white, rounded corners=1.3mm, font=\sffamily\bfseries\normalsize,
          inner sep=0pt, minimum width=8mm, minimum height=8mm] (n) {#2};
    \node[anchor=west, font=\sffamily\bfseries\large, text=cInk] (t) at (n.east) [xshift=3mm] {#3};
  \end{tikzpicture}\par\vspace{1mm}%
  \noindent{\color{#1}\rule{\linewidth}{1pt}}\par\addvspace{2.5mm}%
}
% lighter sub-heading
\newcommand{\wbsub}[2][cBrand]{\par\addvspace{3.5mm}\noindent{\sffamily\bfseries\large\color{#1}#2}\par\addvspace{1.2mm}}
% run-in guided-learning label
\newcommand{\glabel}[2][cBrand]{\par\addvspace{2mm}\noindent{\sffamily\bfseries\color{#1}#2\enspace}\ignorespaces}
% eyebrow tag
\newcommand{\eyebrow}[2][cAccent]{{\sffamily\bfseries\footnotesize\color{#1}\MakeUppercase{#2}}}

% =====================================================================
%  FIGURE / DIAGRAM WRAPPER
% =====================================================================
\newtcolorbox{wbfigure}[1]{
  enhanced, breakable=false, colback=cFigBg, colframe=cRule, boxrule=0.6pt,
  arc=1.6mm, halign=center, left=3mm,right=3mm,top=4mm,bottom=2.6mm,
  after upper={\par\vspace{2.4mm}{\color{cRule}\rule{0.42\linewidth}{0.3pt}}\par\vspace{1.1mm}%
    {\footnotesize\itshape\color{cInk!72}#1}}}

% auto-fit: scale a diagram DOWN to fit the column (with padding margin) if too wide, else leave it.
\newsavebox{\wbscalebox}
\newlength{\wbtargetw}
\newcommand{\wbscale}[1]{%
  \sbox{\wbscalebox}{#1}%
  \setlength{\wbtargetw}{0.94\linewidth}%
  \ifdim\wd\wbscalebox>\wbtargetw
    \resizebox{\wbtargetw}{!}{\usebox{\wbscalebox}}%
  \else
    \usebox{\wbscalebox}%
  \fi}

% TikZ diagram vocabulary
\tikzset{
  wbarrow/.style={-{Latex[length=2.6mm,width=2.1mm]}, line width=0.75pt, draw=cBrand},
  wbarrowg/.style={-{Latex[length=2.6mm,width=2.1mm]}, line width=0.75pt, draw=cInk!45},
  wbbox/.style={draw=cBrand, line width=0.9pt, fill=cBrand!7, rounded corners=1.6mm,
    text=cInk, font=\sffamily\small, align=center, inner sep=3.5pt,
    minimum height=10mm, minimum width=24mm},
  wbboxaccent/.style={wbbox, draw=cAccent, fill=cAccent!12},
  wbboxgood/.style={wbbox, draw=cGood, fill=cGood!10},
  wbboxwarn/.style={wbbox, draw=cSecurity, fill=cSecurity!8},
  wbboxghost/.style={wbbox, draw=cInk!35, fill=cInk!4},
  wbboxdb/.style={wbbox, draw=cDb, fill=cDb!8},
  wbcap/.style={font=\sffamily\scriptsize\itshape, text=cInk!70, align=center},
  wbtag/.style={font=\sffamily\scriptsize, text=cInk!75, align=center},
  wbcyl/.style={cylinder, shape border rotate=90, aspect=0.28, draw=cDb, line width=0.9pt,
    fill=cDb!8, text=cInk, font=\sffamily\small, align=center, minimum height=12mm, minimum width=20mm, inner sep=2pt},
}

% =====================================================================
%  LABS, DEBUGGING, THINKING, REFLECTION, CHECKLISTS
% =====================================================================
% Guided coding lab — step machinery
\newcounter{labstep}
\newcommand{\resetlab}{\setcounter{labstep}{0}}
\newcommand{\labstep}[1]{%
  \par\addvspace{2.5mm}\refstepcounter{labstep}%
  \noindent\begin{tikzpicture}[baseline=(s.base)]
    \node[fill=cLab, text=white, circle, inner sep=0pt, minimum size=6.4mm,
          font=\sffamily\bfseries\footnotesize] (s) {\thelabstep};
    \node[anchor=west, font=\sffamily\bfseries, text=cInk] at (s.east) [xshift=2.2mm] {#1};
  \end{tikzpicture}\par\vspace{0.6mm}%
}
\newcommand{\labwhy}[1]{\noindent{\sffamily\itshape\color{cLab}\faLightbulb\ Why:\ }{\small#1}\par}

% Mini-lab box (Beginner / Intermediate / Production-style)
\newtcolorbox{minilab}[2]{enhanced,breakable,
  colback=cLab!4, colframe=cLab, boxrule=0.7pt, leftrule=3.2pt, arc=1.4mm,
  left=3.2mm,right=3.2mm,top=2.6mm,bottom=2.6mm,
  coltitle=white, colbacktitle=cLab, fonttitle=\sffamily\bfseries\small,
  title={\faFlask\enspace Mini-Lab \textbullet\ #1 \hfill \normalfont\itshape\footnotesize #2}}

% Debugging lab box
\newtcolorbox{debuglab}[1]{enhanced,breakable,
  colback=cDebug!4, colframe=cDebug, boxrule=0.7pt, leftrule=3.2pt, arc=1.4mm,
  left=3.2mm,right=3.2mm,top=2.6mm,bottom=2.6mm,
  coltitle=white, colbacktitle=cDebug, fonttitle=\sffamily\bfseries\small,
  title={\faBug\enspace Debugging Lab \textbullet\ #1}}

% Thinking / reflection prompt box
\newtcolorbox{thinkbox}[1][Think it through]{enhanced,breakable,
  colback=cThink!7, colframe=cThink, boxrule=0.6pt, leftrule=3.2pt, arc=1.4mm,
  left=3.2mm,right=3.2mm,top=2.4mm,bottom=2.4mm,
  coltitle=white, colbacktitle=cThink, fonttitle=\sffamily\bfseries\small,
  title={\faBrain\enspace #1}}

% Build-along project box
\newtcolorbox{buildalong}[1][ShopFlow — extend the capstone]{enhanced,breakable,
  colback=cBuild!5, colframe=cBuild, boxrule=0.8pt, leftrule=3.6pt, arc=1.6mm,
  left=3.4mm,right=3.4mm,top=2.8mm,bottom=2.8mm,
  coltitle=white, colbacktitle=cBuild, fonttitle=\sffamily\bfseries\small,
  title={\faHammer\enspace Build-Along \textbullet\ #1}}

% reflection ruled lines (write-in space)
\newcommand{\reflectionlines}[1][4]{%
  \par\addvspace{1.4mm}%
  \foreach \wbi in {1,...,#1}{%
    \noindent{\color{cRule}\rule{\linewidth}{0.35pt}}\par\vspace{5.4mm}}%
  \vspace{-3mm}\par}

% dotted write-in space with a soft label
\newtcolorbox{writespace}[1][My notes]{enhanced,breakable=false,
  colback=white, colframe=cRule, boxrule=0.5pt, arc=1.2mm,
  left=3mm,right=3mm,top=2mm,bottom=1mm,
  fonttitle=\sffamily\itshape\footnotesize, coltitle=cInk!55, colbacktitle=white,
  title={\faPenNib\ #1}}

% checklists
\newlist{checklist}{itemize}{2}
\setlist[checklist]{label=\textcolor{cBrand}{\large$\square$}, leftmargin=7mm, itemsep=2.2pt, topsep=2pt}
\newlist{tickcheck}{itemize}{1}
\setlist[tickcheck]{label=\textcolor{cGood}{\faCheck}, leftmargin=7mm, itemsep=2pt, topsep=2pt}

% takeaway / plain bullet
\newlist{wbitem}{itemize}{2}
\setlist[wbitem]{label=\textcolor{cAccent}{\footnotesize\faCircle}, leftmargin=6mm, itemsep=2.5pt, topsep=2pt}
\setlist[wbitem,2]{label=\textcolor{cBrand}{\scriptsize\faAngleRight}, leftmargin=5mm}

% =====================================================================
%  CHEAT SHEET
% =====================================================================
\newtcolorbox{cheatsheet}[1][One-Page Cheat Sheet]{enhanced,breakable,
  colback=cBrand!4, colframe=cBrand, boxrule=0.9pt, arc=1.8mm,
  left=3.4mm,right=3.4mm,top=3mm,bottom=3mm,
  coltitle=white, colbacktitle=cBrand, fonttitle=\sffamily\bfseries,
  title={\faClipboardCheck\enspace #1}}
\newcommand{\cheatitem}[2]{\noindent{\sffamily\bfseries\color{cBrandDeep}#1}\ \textemdash\ #2\par\addvspace{0.6mm}}

% =====================================================================
%  INTERVIEW Q / KEYTERM / SUMMARY helpers
% =====================================================================
\newtcolorbox{iqbox}{enhanced,breakable,
  colback=cInterview!5, colframe=cInterview, boxrule=0.7pt, leftrule=3.2pt, arc=1.4mm,
  left=3.2mm,right=3.2mm,top=2.4mm,bottom=2.4mm,
  coltitle=white, colbacktitle=cInterview, fonttitle=\sffamily\bfseries\small,
  title={\faUserTie\enspace Interview questions to master}}
\newcommand{\iq}[1]{\noindent\textcolor{cInterview}{\faAngleRight}\ #1\par\addvspace{0.8mm}}

\newcommand{\keyterm}[1]{\textbf{\color{cBrandDeep}#1}}

% two-col best/mistakes side by side helper
\newenvironment{twocolcallouts}{\par\begin{minipage}[t]{0.485\linewidth}}{\end{minipage}\hfill\ignorespaces}
\newcommand{\colgap}{\end{minipage}\hfill\begin{minipage}[t]{0.485\linewidth}}

% =====================================================================
%  TABLES
% =====================================================================
\newcommand{\wbtablehead}[1]{\cellcolor{cBrand}\textcolor{white}{\sffamily\bfseries #1}}
\newcolumntype{L}[1]{>{\raggedright\arraybackslash}p{#1}}
\newcolumntype{C}[1]{>{\centering\arraybackslash}p{#1}}
\definecolor{cRowA}{HTML}{F3F5FA}
\newcommand{\wbrowcolors}{\rowcolors{2}{white}{cRowA}}

% =====================================================================
%  CHAPTER OPENER
% =====================================================================
\usepackage{fancyhdr}
\usepackage{titlesec}
\usepackage[hidelinks]{hyperref}
\hypersetup{bookmarksnumbered=true, pdfstartview=FitH,
  colorlinks=true, linkcolor=cBrand, urlcolor=cProd, citecolor=cBrand}

\newcommand{\daychapter}[5]{%
  % #1 day number (01..14)  #2 track eyebrow  #3 title  #4 subtitle  #5 study time
  \clearpage
  \phantomsection
  \addcontentsline{toc}{chapter}{\texorpdfstring{\textbf{Day #1}\, — #3}{Day #1 — #3}}%
  \markboth{Day #1 \ \textbullet\ \ #3}{}%
  \begin{tcolorbox}[enhanced, breakable=false, colback=cBrand, colframe=cBrandDeep,
      boxrule=0pt, arc=2.2mm, left=6mm,right=6mm,top=5mm,bottom=5.5mm,
      overlay={\node[anchor=north east, font=\sffamily\bfseries, text=white!12,
         scale=5.4] at ([xshift=-3mm,yshift=2mm]frame.north east) {#1};}]
    {\sffamily\bfseries\footnotesize\textcolor{cAccent}{\MakeUppercase{Day #1 \ \ \textbullet\ \ #2}}}\\[1.6mm]
    {\sffamily\bfseries\huge\textcolor{white}{#3}}\\[2.2mm]
    {\sffamily\normalsize\textcolor{white!88}{#4}}\\[3.4mm]
    {\sffamily\footnotesize\textcolor{white}{\faRegularClock\ \ Estimated study time: \textbf{#5}}}
  \end{tcolorbox}\par\vspace{3mm}%
}
% robust clock (regular style may vary): alias
\providecommand{\faRegularClock}{\faClock}

% learning objectives box
\newtcolorbox{objectives}{enhanced,breakable,
  colback=cAccent!7, colframe=cAccent, boxrule=0.7pt, leftrule=3.4pt, arc=1.4mm,
  left=3.2mm,right=3.2mm,top=2.6mm,bottom=2.6mm,
  coltitle=cInk, colbacktitle=cAccent!22, fonttitle=\sffamily\bfseries\small,
  title={\faBullseye\enspace By the end of this day, you will be able to\ldots}}

% =====================================================================
%  HEADERS / FOOTERS
% =====================================================================
\pagestyle{fancy}
\fancyhf{}
\fancyhead[L]{\footnotesize\sffamily\color{cInk!58}\nouppercase{\leftmark}}
\fancyhead[R]{\footnotesize\sffamily\color{cInk!58}The Backend Engineer's Lab Notebook}
\fancyfoot[C]{\footnotesize\sffamily\color{cInk!55}\thepage}
\renewcommand{\headrulewidth}{0pt}
\renewcommand{\footrulewidth}{0pt}
\fancyhead[L]{\footnotesize\sffamily\color{cInk!58}\nouppercase{\leftmark}%
  \protect\raisebox{-3mm}{\protect\makebox[0pt][l]{\protect\color{cRule}\protect\rule{\headwidth}{0.4pt}}}}

\fancypagestyle{plain}{\fancyhf{}\fancyfoot[C]{\footnotesize\sffamily\color{cInk!55}\thepage}\renewcommand{\headrulewidth}{0pt}}

% paragraph spacing for readability
\setlength{\parindent}{0pt}
\setlength{\parskip}{2.6pt plus 1pt}
\linespread{1.045}

% line-breaking safeguards (prevents long inline-code tokens overflowing the margin)
\tolerance=1400
\hbadness=2500
\emergencystretch=2.4em
\hyphenpenalty=200
\pretolerance=1200

% spacing niceties
\setlist{topsep=2pt,partopsep=0pt}
\raggedbottom
